\subsection{向量丛与主纤维化}

7.10.1. 设 F 是一个巴拿赫空间。若向量丛 M 的所有纤维 \(M_{b}\)（\(b \in B\)）都与 F（作为巴拿赫空间）同构，则称以 B 为基底的向量丛 M 是 F 型纯向量丛。

设 M 是以 B 为基底的 F 型纯向量丛，令 P 为向量丛 \(\mathcal{L}(\mathrm{F}_{\mathrm{B}}; \mathrm{M})\) 的开子流形，它由满足以下条件的对 \((b, u)\) 组成：\(b \in \mathbf{B}\)，且 \(u\) 是从 \(\mathrm{F}_{b} = \mathrm{F}\) 到 \(\mathbf{M}_{b}\) 的同构。F 的自同构群 \(\mathbf{GL}(\mathbf{F})\) 在 P 上右作用，规定 \((b, u). g = (b, u \circ g)\)，其中 \((b, u) \in \mathbf{P}\) 且 \(g \in \mathbf{GL}(\mathbf{F})\)。将 \(\pi_{\mathbf{P}}\) 记为从 P 到 B 的映射 \((b, u) \mapsto b\)。四元组 \(\lambda = (\mathbf{P}, \mathbf{GL}(\mathbf{F}), \mathbf{B}, \pi_{\mathbf{P}})\)（其中 \(\mathbf{GL}(\mathbf{F})\) 赋有其作为群流形的典范结构（5.12.2））是一个主纤维化（6.2.1），称为 M 的标架纤维化。映射 \(((b, u), h) \mapsto u(h)\)，从 \(\mathbf{P} \times \mathbf{F}\) 到 M，赋予 M 一个与 \(\lambda\) 相伴、典型纤维为 F 的纤维丛结构（6.5.1）。

当 \(\mathbf{F} = \mathbf{K}^n\) 时，可以把一个同构 \(u\)，从 \(\mathbf{F}\) 到 \(\mathbf{M}_b\)，与 \(\mathbf{M}_b\) 的一个基对应起来：该基由 \(u\) 作用于 \(\mathbf{K}^n\) 的典范基所得的像组成。于是，\(\mathbf{M}\) 的标架丛可与 \(\mathbf{M} \times_{\mathbf{B}} \ldots \times_{\mathbf{B}} \mathbf{M}\) 的开子流形认同，该子流形由不同纤维的基 \((e_1, \ldots, e_n)\)，其中这些纤维是 \(\mathbf{M}_b\)，组成。

设 U 是 B 的一个开子集，\(t = (\mathrm{U}, \varphi, \mathrm{E})\) 是定义域为 U 且 E = F 的 M 的一个向量丛图。于是映射 \(b \mapsto (b, t_{b})\) 是 P 的一个截面，记为 \(\tilde{t}\)；映射 \(t \mapsto \tilde{t}\) 是从形如 \((\mathrm{U}, \varphi, \mathrm{F})\) 的 M 的向量丛图集到 P 在 U 上的截面集的双射。

7.10.2. 反之，设 \(\lambda = (\mathbf{Q}, \mathbf{G}, \mathbf{B}, \pi_{\mathbf{Q}})\) 是一个主纤维化，\(\varphi\) 是从群流形 G 到巴拿赫空间 F 的自同构群 \(\mathbf{GL}(\mathbf{F})\) 的群流形同态。令 G 在 F 上左作用，规定 \(g \cdot h = \varphi(g)(h) (h \in \mathbf{F}, g \in \mathbf{G})\)。设 M 是与 \(\lambda\) 相伴且典型纤维为 F 的纤维丛，\(\rho: \mathbf{Q} \times \mathbf{F} \to \mathbf{M}\) 是其标架映射（6.5.1）。令 \(\pi\) 为从 M 到 B 的映射，由下式定义：

\[
\pi (\rho (q, h)) = \pi_ {\mathrm{Q}} (q) \quad (q \in \mathbf{Q}, h \in \mathbf{F}).
\]

设 s 是 Q 在 B 的一个开集 U 上的截面；映射

\[
\tilde {s}: (b, h) \mapsto (s (b), h)
\]

于是它是从 U \(\times\) F 到 \(\pi^{-1}(\mathbf{U})\) 的双射。在对 (M，\(\pi\)) 上存在一个以 B 为基底的向量丛结构，并且只有一个这样的结构使得三元组 \(t_{s} = (\mathbf{U}, \tilde{s}^{-1}, \mathbf{F})\) 是向量丛图（对 Q 的每个截面 s 都如此）。这个结构的底层流形结构就是与 \(\lambda\) 相伴的纤维丛的结构。

设 \(q \in \mathbf{Q}\)；置 \(b = \pi_{\mathbf{Q}}(q)\)，并令 \(u\) 为从 \(\mathbf{F}\) 到 \(\mathbf{M}_b\) 的同构，它满足 \(u(h) = \rho(q, h)\)，其中 \(h \in \mathbf{F}\)。映射 \(f: q \mapsto (b, u)\) 是一个与 \(\varphi\) 相容的主纤维化 B-态射，从 \((\mathbf{Q}, \mathbf{G}, \mathbf{B}, \pi_{\mathbf{Q}})\) 到标架丛 \((\mathbf{P}, \mathbf{GL}(\mathbf{F}), \mathbf{B}, \pi_{\mathbf{P}})\)，它是向量丛 \(\mathbf{M}\) 的标架丛。

7.10.3. 下面沿用 7.6 节的记号。对每个 \(i \in I\)，令

\[
\lambda_ {i} = (\mathbf {P} _ {i}, \mathbf {G} _ {i}, \mathbf {B}, \pi_ {i})
\]

这是一个以 B 为基底的主纤维化，并假设 \(G_{i}\) 通过一个同态在巴拿赫空间 \(V_{i}\) 上左作用

\[
\varphi_{i}: \mathbf{G}_{i} \rightarrow \mathbf{GL}(\mathrm{V}_{i}).
\]

设 \(\mathbf{M}_i\) 是与 \(\lambda_i\) 相伴且典型纤维为 \(\mathbf{V}_i\) 的纤维丛。置 \(\mathcal{M} = (\mathbf{M}_i)_{i\in I}\) 和 \(\mathcal{V} = (\mathrm{V}_i)_{i\in \mathbf{I}}\)，并令 \(\lambda\) 为各 \(\lambda_{i}\) 在 B 上的主积纤维化（6.2.5）。置 \(\lambda = (\mathrm{P},\mathrm{G},\mathrm{B},\pi_{\mathbf{P}})\)，其中 \(\mathbf{G} = \prod_{i\in \mathbf{I}}\mathbf{G}_i\)。

现在设 \(\tau\) 是一个向量函子。对于 \(g = (g_i) \in G\)，令 \(\varphi(g)\) 为 \(\operatorname{Hom}(\mathcal{V}, \mathcal{V})\) 中由下式定义的元素：

\[
\begin{array}{l l} \varphi (g) _ {i} = \varphi_ {i} (g _ {i}) & \text { if } \quad i \in \mathrm{I} _ {+} \\ \varphi (g) _ {i} = \varphi_ {i} (g _ {i}) ^ {- 1} & \text { if } \quad i \in \mathrm{I} _ {-}. \end{array}
\]

于是群 G 在 \(\tau(\mathcal{V})\) 上通过态射 \(g\mapsto\tau(\varphi(g))\)（从 G 到 \(\mathbf{GL}(\tau(\mathcal{V}))\)）作用。

此外，设 \(x = (x_i)\) 是 P 的一个点，且 \(b = \pi_{\mathbf{P}}(x)\)。对每个 \(i\)，第 6.5.2 号中定义的映射 \(\theta_{x_i}\) 是从 \(V_i\) 到 \((\mathbf{M}_i)_b\) 的同构。令 \(\theta_x\) 为 \(\operatorname{Hom}(\mathcal{V}, ((\mathbf{M}_i)_b)_{i \in I})\) 中由下式定义的元素：

\[
\begin{array}{l l} (\theta_ {x}) _ {i} = \theta_ {x _ {i}} & \text { if } \quad i \in I _ {+} \\ (\theta_ {x}) _ {i} = \theta_ {x _ {i}} ^ {- 1} & \text { if } \quad i \in I _ {-}. \end{array}
\]

令 \(\rho\) 为映射 \((x, h) \mapsto (b, \tau(\theta_x)(h))\)，从 \(\mathbf{P} \times \tau(\mathcal{V})\) 到向量丛 \(\tau(\mathcal{M})\)；映射 \(\rho\) 赋予 \(\tau(\mathcal{M})\) 一个与 \(\lambda\) 相伴、典型纤维为 \(\tau(\mathcal{V})\) 的纤维丛结构。

这些考虑可推广到有限维向量函子的情形，或推广到带有有限维 K-代数结构的域 L 上的向量函子的情形。

